\section{Introduction}

The growth of index investing creates a simple expectation: when more capital follows an index, changes in its membership should generate larger purchases and sales. Yet \citet{greenwood2025} document a striking development in the United States. Average returns around additions to the S\&P 500 fell from 7.4 \% in the 1990s to less than 1 \% in their latest decade, despite the expansion of index investing. Their study asks how markets have become able to accommodate these trades with a smaller price response. This thesis takes that question to Sweden.

We ask whether the disappearing index effect also appears around changes in the OMX Stockholm Benchmark Index (OMXSB), and what reported fund trading and the subsequent return pattern reveal about the association that remains. Our main finding is that later-decade additions continue to earn positive average market-adjusted returns. Their mean is 2.15 \% during 2017--2026. The historical pattern is nevertheless uneven, and the latest five-year estimate is less precise. Sweden therefore provides evidence of a positive later-period addition association without establishing that the effect is stable, permanent or caused by index membership.

The setting allows us to connect this return comparison to two features of the Swedish market. First, OMXSB changes occur in scheduled reviews, with several securities sharing an announcement and implementation date. The index selects securities using trading activity and industry representation, rather than a single market-capitalization boundary. Both the composition of additions and the time between announcement and implementation can consequently matter for measured returns. Second, Swedish companies often have more than one listed share class. When one class changes index status while another retains its status, their relative performance provides an additional perspective on whether the return movement is concentrated in the affected security or shared within the firm.

Our event register contains 536 changes across 39 reviews from July 2007 through June 2026. After return-data requirements and the stated event exclusions, the principal sample contains 256 additions and 215 deletions. Following the anchor study's event-window approach, we examine announcement returns, implementation returns and the interval from the day before announcement through the day after implementation. Our primary cumulative abnormal return is the sum of daily stock returns less the OMXSB gross-index return. We combine these returns with dated fund classifications and reported holdings, and compare the commercial holdings measure with regulatory fund reports. Inference clusters events by review because their dates and market conditions coincide.

Addition returns are substantially lower at the end of the sample than at its beginning. The four period means are 8.52, 0.25, 2.85 and 1.50 \%. This is a decline between endpoints, not a steady downward trajectory. Pooling 2017--2026 gives a mean of 2.15 \% and a 95 \% review-bootstrap interval of 1.01 to 3.39 \%. The result remains positive under the reported extreme-return and event-exclusion sensitivities. By contrast, the latest period alone has an interval extending from approximately -0.10 to 3.22 \% and is sensitive to an unusually large addition return. Later-decade deletions average -1.04 \%, with an interval spanning zero. The evidence for additions is therefore clearer in the pooled later period than in the latest block, and clearer than the corresponding deletion evidence.

The holdings data show that these returns occur alongside economically meaningful benchmark-linked trading. In the common later-decade addition sample, covered funds buy an average of 1.03 \% of shares outstanding, while mean CAR is 2.71 \%. Their ratio is 2.63 percentage points of return per percentage point of measured purchases. This gives the multiplier used by Greenwood and Sammon an intuitive Swedish counterpart. It is a ratio of observed averages, however, rather than an estimate of the causal response to an additional purchase. The reporting interval is broader than the return window, coverage changes over time, and unobserved investors may trade in either direction. A matched comparison with Finansinspektionen reports supports a similar later-period demand scale without resolving those identification problems.

We then examine the explanations emphasized by the anchor study. Adjusting for observable firm characteristics does not eliminate the large earliest-to-latest difference. Pre-announcement returns are substantial, but their period pattern does not show a steady increase matching the decline in the full event-window return. Quoted spreads narrow and trading becomes more concentrated around implementation, consistent with changes in how reconstitutions are traded. These observations help describe the market in which the return pattern occurs. They do not isolate improved liquidity provision as its cause. Nor does our concurrent-index check reconstruct the offsetting demand from migrations between the S\&P 500 and S\&P MidCap 400. The Swedish and US estimates remain institutionally and statistically distinct comparisons.

A further result concerns the timing of the measurement. The announcement-to-implementation gap becomes shorter during the sample. On identical events, we therefore compare the full return window with a common initial window from AD-1 through AD+5, ending before implementation in every review. Initial-window addition means are 1.19 \% in the earliest block and 1.80 \% in the latest. The observed full-window decline is concentrated in the remaining interval. This exploratory decomposition does not establish equal initial responses or show that shorter notice causes lower returns. It demonstrates why changing the duration of the event window matters when interpreting historical changes in the measured index effect.

Two further comparisons help describe the remaining return. After correcting a one-day overlap between its pre-return control and outcome, a size-screened comparison gives a later addition coefficient of 1.65 percentage points, with an interval of [0.48, 2.83]. On a fixed later-decade event cohort, cumulative CAR through ED+20 is 3.30\%, with an interval of [1.31\%, 5.28\%]. The observed mean has not returned to zero by this finite horizon. Neither exercise supplies the missing counterfactual for a causal interpretation.

The share-class evidence adds another constraint on interpretation. Across 39 selected events, sign-adjusted log market-excess returns average 3.18 \% for the class changing status and 2.85 \% for its sister class. The relative mean is 0.33 percentage points and is imprecisely estimated. Much of the observed movement in this subset is thus shared by the two classes. Shared firm information could produce such a pattern, but so could trading that transmits demand across related securities. Differences in voting rights, liquidity and investors also prevent treating the sister class as a perfect counterfactual. The result cannot establish a causal upper bound on index demand, and it is not a decomposition of the full return sample.

Our contribution is a replication and extension that connects the historical return pattern to its trading and measurement counterparts. The later Swedish addition association remains positive alongside reported fund buying and concentrated implementation trading. The timing comparison shows how an institutional calendar change affects the interpretation of the historical return contrast, while the holdings and share-class exercises clarify what the measured association can represent. Together, these results extend the anchor study's questions without converting an observational event study into a structural demand estimate. Section 2 reviews the closest literature, Section 3 presents the data and methodology, Section 4 reports the evidence, and Sections 5 and 6 discuss the findings and conclude.
